\section{The Jacobian action}
\label{sec:jac}

The solver applies $\partial R/\partial q$ to a direction $v$ without assembling it. Two variants
exist and they differ in whether the deferred correction and the Rhie--Chow reconstruction are
differentiated.

\subsection{Lagged and exact}

The \emph{lagged} action differentiates the first-order operator only: the correction
$\mathbf{h}_\scs$ contributes nothing, and the Rhie--Chow term is differentiated with the nodal
pressure gradient held fixed. This is the operator the assembled matrix can represent, because
including either term widens the pressure coupling past nearest neighbours and breaks the
one-ring sparsity.

The \emph{exact} action differentiates both. The costs are asymmetric and worth separating:

\begin{itemize}
\item Differentiating the Rhie--Chow correction, Eq.~\eqref{eq:corr}, requires the
      \emph{direction's} nodal pressure gradient, so a reconstruction pass runs per matvec. Its
      absence is not visible in the residual and does not make any single Newton step fail --- it
      caps Newton at a linear rate, and only where the continuity rows matter. On a
      backward-facing step, omitting it put the Jacobian $11.7\%$ away from a finite difference of
      the residual in the continuity rows while momentum stayed at $0.3\%$.
\item Differentiating the correction requires the direction's nodal \emph{velocity} gradient ---
      nine components against the three the pressure term needs --- plus the limiter's own
      derivative.
\end{itemize}

Both direction gradients are produced by one reconstruction sweep rather than two, since the pass
is parameterised by field count (\code{NC}).

\subsection{The coefficient's velocity sensitivity}

$G_\scs$ in Eq.~\eqref{eq:coeff} depends on the state through $|\bar{\uvec}|^2$, so the Jacobian
must carry $\partial G/\partial\uvec$. Writing $G = K/\sqrt{S}$ with
$K = \alpha|\Avec|^2/(\Avec\cdot\dvec)$ and $S$ the bracket of Eq.~\eqref{eq:coeff}, the
derivative along $v$ with face average $\bar{v}$ is
\begin{equation}
\delta G \;=\; -\,g\,\bigl(\bar{\uvec}\cdot\bar{v}\bigr),
\qquad
g \;=\; \frac{4\,u^2_{\text{scale}}}{h^2}\,\frac{G}{S} \;=\; G^3\,w ,
\end{equation}
with $w$ exactly Eq.~\eqref{eq:wdu}. The $G^3 w$ factorisation is the one used because $w$ is pure
geometry: it can be tabulated with the mesh, leaving the surface loop one multiply. The
Rhie--Chow mass flux $-G c$ therefore gains $g\,c\,(\bar{\uvec}\cdot\bar{v})$.

\subsection{Limiter derivatives and subgradients}
\label{sec:jac:subgrad}

The exact action needs $\partial\Lambda/\partial\Delta$, and two of the three limiters are not
differentiable everywhere. How that is handled is a correctness matter, not a detail.

\paragraph{The clip} is discontinuous in its derivative at both bounds. Naively differentiating
each branch makes the result depend on which side of the bound an ulp of round-off lands, and
since the two sides differ by a finite amount, an ulp flips an $O(1)$ quantity. In practice this
made the Jacobian action \emph{layout dependent}: the packed and standard arms disagreed at
$3\times10^{-3}$, because FMA contraction differs between them and tipped branches at ties. The
fix is a rounding band: within a tolerance scaled by the magnitudes involved, the derivative takes
the value that does not depend on which side the comparison fell.

\paragraph{Darwish--Moukalled} has the same defect at $a=0$ and $b=0$, where the absolute values
have no derivative. The two one-sided limits differ by $2\,\delta a\,b/\mathrm{den}$ and
$2\,a\,\delta b/\mathrm{den}$ respectively. Inside a rounding band the implementation takes the
\emph{minimum-norm subgradient}, which is zero --- the symmetric element of the subdifferential of
$|\cdot|$ at zero. Neither side is preferred, and it is the only choice independent of where an
ulp lands.

\paragraph{Venkatakrishnan} branches too, on $\Delta \ge 0$, but needs no band: at $\Delta = 0$ its
$\psi$ is exactly $1$ and every term carrying the branch is multiplied by $\Delta$, so the
discontinuity cancels and the derivative is continuous there. This was checked rather than assumed.

\subsection{Partial assembly}

A third variant precomputes the surface quantities that do not depend on the direction and stores
them, so the apply reads a compressed tangent instead of the state
(\code{cvfem\_hex8\_build\_pa\_tangent}). Per surface it stores the upwind split, and for the
Rhie--Chow arms the combined velocity-sensitivity factor $g\,c\,\bar{\uvec}$ as three components.
The store is sixty scalars per element.

This is where the economics of precomputation invert relative to the residual. The Jacobian is
applied many times at a \emph{fixed} state, so the build amortises over the Krylov iterations; the
residual is evaluated once per state, so precomputing for it would mean rebuilding every call.
That asymmetry is why the Jacobian reads a tabulated $G_\scs$ while the residual recomputes it
inline, and why the same change --- taking $\dvec$ from the Jacobian --- helped the residual
substantially and the Jacobian not at all until $w$ was tabulated too.

\subsection{Verification}

Three oracles, each catching a different class of error:

\begin{itemize}
\item \emph{Cross-layout}: the packed and standard actions compared node by node, which catches
      staging and scatter errors. Reads $\sim10^{-16}$.
\item \emph{Finite difference of the residual}: catches a missing or wrong term in the
      linearisation. The exact arm is gated; the lagged arm is gated \emph{from below}, because a
      lagged action that agreed with the finite difference would mean the correction never reached
      the residual and the comparison is vacuous.
\item \emph{Partially assembled against direct}: catches disagreement between a producer of
      precomputed data and its consumer. This is the only one of the three that runs on a
      non-affine mesh, and consequently the only one that can see a geometry inconsistency ---
      Section~\ref{sec:geom:edges} records it catching one at $1.4\times10^{-5}$ that the other
      two could not.
\end{itemize}

A verification must include whichever passes the timed apply included. Comparing an action that
ran the boundary closure against a reference that did not reports the closure as a staging error,
which has happened and reads as an $O(1)$ failure.
